% GENERATED FILE. Edit canonical lessons, then run npm run generate:books.
% canonical-source-hash: fnv1a64:629a15b9378ad564
% canonical-lessons: TE-C105-stesan, TE-C105-bassu, TE-C105-samarthyam, TE-C105-naipunyam, TE-C105-hakku

\chapter{Station, Bus, Ability, Skill, Right}
\label{ch:te-station-bus-ability-skill-right}
\hlchaptermodality{105}

\begin{chapteropening}
\textbf{By the end of this chapter:} \emph{I can say a station, a bus, ability, a skill and a right.}

Complete the last lesson of chapter 105: i can say a station, a bus, ability, a skill and a right.
\end{chapteropening}

\section[\texorpdfstring{\te{స్టేషన్} (sṭēṣan)}{స్టేషన్ (sṭēṣan)}]{\te{స్టేషన్} (sṭēṣan) — a station}

\label{lesson:TE-C105-stesan}



\begin{quote}
Before the new one: say the Telugu for a corner, then the Telugu for a turn.
\end{quote}

\subsection*{You'll want to know: \te{స్టేషన్}}
\textbf{\te{స్టేషన్}} — \emph{sṭēṣan} — \textquotedblleft{}a station\textquotedblright{}.

\begin{grammarlens}[title={What you've built}]
One more word for this chapter.
\end{grammarlens}

\subsection*{Your turn}
\begin{itemize}
\raggedright
  \item \emph{Say it:} \emph{sṭēṣan}
  \item \emph{Say it:} \emph{sṭēṣan}, once more
  \item \emph{Say it:} say \emph{sṭēṣan}
  \item \emph{Recall:} say the Telugu for a corner, then the Telugu for a turn, then say \emph{sṭēṣan} again
\end{itemize}

\begin{tcolorbox}[breakable,colback=teal!4,colframe=teal!35!black,title={Before you move on}]
What does \te{స్టేషన్} mean? (\textquotedblleft{}A station\textquotedblright{}.) Say it once more.
\end{tcolorbox}
\section[\texorpdfstring{\te{బస్సు} (bassu)}{బస్సు (bassu)}]{\te{బస్సు} (bassu) — a bus}

\label{lesson:TE-C105-bassu}



\begin{quote}
Before the new one: say the Telugu for a turn, then the Telugu for a station.
\end{quote}

\subsection*{You'll want to know: \te{బస్సు}}
\textbf{\te{బస్సు}} — \emph{bassu} — \textquotedblleft{}a bus\textquotedblright{}.

\begin{grammarlens}[title={What you've built}]
One more word for this chapter.
\end{grammarlens}

\subsection*{Your turn}
\begin{itemize}
\raggedright
  \item \emph{Say it:} \emph{bassu}
  \item \emph{Say it:} \emph{bassu}, once more
  \item \emph{Say it:} say \emph{bassu}
  \item \emph{Recall:} say the Telugu for a turn, then the Telugu for a station, then say \emph{bassu} again
\end{itemize}

\begin{tcolorbox}[breakable,colback=teal!4,colframe=teal!35!black,title={Before you move on}]
What does \te{బస్సు} mean? (\textquotedblleft{}A bus\textquotedblright{}.) Say it once more.
\end{tcolorbox}
\section[\texorpdfstring{\te{సామర్థ్యం} (sāmarthyaṁ)}{సామర్థ్యం (sāmarthyaṁ)}]{\te{సామర్థ్యం} (sāmarthyaṁ) — ability}

\label{lesson:TE-C105-samarthyam}



\begin{quote}
Before the new one: say the Telugu for a station, then the Telugu for a bus.
\end{quote}

\subsection*{You'll want to know: \te{సామర్థ్యం}}
\textbf{\te{సామర్థ్యం}} — \emph{sāmarthyaṁ} — \textquotedblleft{}ability\textquotedblright{}.

\begin{grammarlens}[title={What you've built}]
One more word for this chapter.
\end{grammarlens}

\subsection*{Your turn}
\begin{itemize}
\raggedright
  \item \emph{Say it:} \emph{sāmarthyaṁ}
  \item \emph{Say it:} \emph{sāmarthyaṁ}, once more
  \item \emph{Say it:} say \emph{sāmarthyaṁ}
  \item \emph{Recall:} say the Telugu for a station, then the Telugu for a bus, then say \emph{sāmarthyaṁ} again
\end{itemize}

\begin{tcolorbox}[breakable,colback=teal!4,colframe=teal!35!black,title={Before you move on}]
What does \te{సామర్థ్యం} mean? (\textquotedblleft{}Ability\textquotedblright{}.) Say it once more.
\end{tcolorbox}
\section[\texorpdfstring{\te{నైపుణ్యం} (naipuṇyaṁ)}{నైపుణ్యం (naipuṇyaṁ)}]{\te{నైపుణ్యం} (naipuṇyaṁ) — a skill}

\label{lesson:TE-C105-naipunyam}



\begin{quote}
Before the new one: say the Telugu for a bus, then the Telugu for ability.
\end{quote}

\subsection*{You'll want to know: \te{నైపుణ్యం}}
\textbf{\te{నైపుణ్యం}} — \emph{naipuṇyaṁ} — \textquotedblleft{}a skill\textquotedblright{}.

\begin{grammarlens}[title={What you've built}]
One more word for this chapter.
\end{grammarlens}

\subsection*{Your turn}
\begin{itemize}
\raggedright
  \item \emph{Say it:} \emph{naipuṇyaṁ}
  \item \emph{Say it:} \emph{naipuṇyaṁ}, once more
  \item \emph{Say it:} say \emph{naipuṇyaṁ}
  \item \emph{Recall:} say the Telugu for a bus, then the Telugu for ability, then say \emph{naipuṇyaṁ} again
\end{itemize}

\begin{tcolorbox}[breakable,colback=teal!4,colframe=teal!35!black,title={Before you move on}]
What does \te{నైపుణ్యం} mean? (\textquotedblleft{}A skill\textquotedblright{}.) Say it once more.
\end{tcolorbox}
\section[\texorpdfstring{\te{హక్కు} (hakku)}{హక్కు (hakku)}]{\te{హక్కు} (hakku) — a right}

\label{lesson:TE-C105-hakku}



\begin{quote}
Before the new one: say the Telugu for ability, then the Telugu for a skill.
\end{quote}

\subsection*{You'll want to know: \te{హక్కు}}
\textbf{\te{హక్కు}} — \emph{hakku} — \textquotedblleft{}a right\textquotedblright{}.

\begin{grammarlens}[title={What you've built}]
One more word for this chapter.
\end{grammarlens}

\subsection*{Your turn}
\begin{itemize}
\raggedright
  \item \emph{Say it:} \emph{hakku}
  \item \emph{Say it:} \emph{hakku}, once more
  \item \emph{Say it:} say \emph{hakku}
  \item \emph{Recall:} say the Telugu for ability, then the Telugu for a skill, then say \emph{hakku} again
\end{itemize}

\begin{tcolorbox}[breakable,colback=teal!4,colframe=teal!35!black,title={Before you move on}]
What does \te{హక్కు} mean? (\textquotedblleft{}A right\textquotedblright{}.) Say it once more.
\end{tcolorbox}
